% ======================================================================
% Diagrama de flujo — generado automáticamente por generar_diagrama_flujo.py
% NO EDITAR MANUALMENTE; regenerar desde el descriptor JSON.
% ======================================================================
\documentclass[11pt,a4paper]{article}

% ── Codificación, fuentes y lenguaje ───────────────────────────────────────────
\usepackage[utf8]{inputenc}
\usepackage[T1]{fontenc}
\usepackage[default,scale=0.95]{sourcesanspro}
\renewcommand{\familydefault}{\sfdefault}
\usepackage[spanish,es-noshorthands,es-tabla]{babel}

% ── Geometría y espaciado ─────────────────────────────────────────────────────
\usepackage[top=2.2cm, bottom=2.5cm, left=2.2cm, right=2.2cm, headheight=15pt]{geometry}
\usepackage{setspace}
\setstretch{1.18}
\usepackage{parskip}

% ── TikZ (simbolos ISO 5807 / ANSI X3.5) ──────────────────────────────────────
\usepackage{tikz}
\usetikzlibrary{babel, arrows.meta, positioning, shapes.geometric,
                shapes.symbols, decorations.pathmorphing, calc}

% ── Tablas y listas ───────────────────────────────────────────────────────────
\usepackage{booktabs}
\usepackage{tabularx}
\usepackage{array}
\usepackage{enumitem}

% ── Colores, cajas e iconos ───────────────────────────────────────────────────
\usepackage[dvipsnames]{xcolor}
\usepackage{tcolorbox}
\tcbuselibrary{skins, breakable}
\usepackage{fontawesome5}

% ── Secciones con barra de color (infografía) ─────────────────────────────────
\usepackage{titlesec}
\titleformat{\section}
  {\normalfont\LARGE\bfseries\color{azulNoche}}
  {\colorbox{cyanNeon}{\color{fondoOscuro}\thesection}}{0.7em}{}
  [\vspace{2pt}\color{cyanNeon}\rule{\linewidth}{1.8pt}]
\titlespacing*{\section}{0pt}{24pt}{10pt}
\titleformat{\subsection}
  {\normalfont\large\bfseries\color{azulNoche}}
  {\textcolor{cyanNeon}{\thesubsection}}{0.7em}{}
  [\vspace{1pt}\color{grisLinea}\rule{\linewidth}{0.6pt}]
\titlespacing*{\subsection}{0pt}{14pt}{6pt}

% ── Cabeceras y pies de página ────────────────────────────────────────────────
\usepackage{fancyhdr}
\usepackage{lastpage}
\pagestyle{fancy}
\fancyhf{}
\renewcommand{\headrulewidth}{0pt}
\renewcommand{\footrulewidth}{0pt}
\fancyfoot[C]{%
  \begin{minipage}{\linewidth}
    \vspace{2pt}
    \color{cyanNeon}\rule{\linewidth}{1.2pt}\\[2pt]
    \centering\color{grisTexto}\footnotesize
    Página \thepage\ de \pageref{LastPage}
  \end{minipage}%
}

% ── Hiperenlaces ──────────────────────────────────────────────────────────────
\usepackage[colorlinks=true, linkcolor=azulNoche, urlcolor=cyanNeon]{hyperref}
\sloppy
\emergencystretch 3em

% ── Paleta de colores ─────────────────────────────────────────────────────────
\definecolor{fondoOscuro}{HTML}{0D1B2A}
\definecolor{verdeTurquesa}{HTML}{004D40}
\definecolor{azulNoche}{HTML}{1B3A6B}
\definecolor{azulMedio}{HTML}{2A5298}
\definecolor{cyanNeon}{HTML}{00C8E0}
\definecolor{verdeSignal}{HTML}{00B887}
\definecolor{naranjaVivo}{HTML}{FF7043}
\definecolor{rojoAlerta}{HTML}{E53935}
\definecolor{grisPapel}{HTML}{F4F6FA}
\definecolor{grisLinea}{HTML}{C8D0E0}
\definecolor{grisTexto}{HTML}{6B7A99}
\definecolor{amarilloNota}{HTML}{FFB300}

% ── Tarjetas ──────────────────────────────────────────────────────────────────
\tcbset{
  cajaContenido/.style={
    enhanced, breakable,
    arc=5pt, outer arc=5pt,
    colback=grisPapel, colframe=grisLinea,
    boxrule=0.8pt,
    fonttitle=\bfseries\small, coltitle=azulNoche,
    top=6pt, bottom=6pt, left=10pt, right=10pt,
  },
}

% ── Macros ────────────────────────────────────────────────────────────────────
\newcommand{\iconotexto}[2]{\textcolor{cyanNeon}{\faIcon{#1}}~#2}
\newcommand{\iconoBanda}{sitemap}
\newcommand{\bandaTitulo}[3][fondoOscuro]{%
  \noindent
  \begin{tcolorbox}[%
    enhanced, arc=4mm, colback=#1!10!white, colframe=#1, boxrule=0pt,
    borderline west={6pt}{0pt}{cyanNeon},
    top=6pt, bottom=6pt, left=8pt, right=8pt,
  ]
    \noindent
    \begin{minipage}[c]{0.10\linewidth}
      \centering\color{cyanNeon}\faIcon{exchange-alt}
    \end{minipage}%
    \hspace{8pt}%
    \begin{minipage}[c]{0.88\linewidth}
      {\fontsize{20}{24}\selectfont\bfseries\color{white}#2}\par\vspace{5pt}%
      \textcolor{cyanNeon!80!white}{\small\faIcon{angle-right}~#3}%
    \end{minipage}
    \vspace{4pt}
  \end{tcolorbox}%
  \vspace{8pt}%
}


  % ── Estilos de flujo (ISO 5807 / ANSI X3.5) ─────────────────────────────
  \tikzset{
    flujo/.style={
      >=Stealth,
      font=\footnotesize\sffamily,
    },
    terminador/.style={
      draw=azulNoche, fill=azulNoche!8!white, rounded corners=3pt,
      text width=1.6cm, align=center, inner sep=2mm,
      font=\footnotesize\sffamily,
    },
    proceso/.style={
      draw=azulNoche, fill=cyanNeon!10!white,
      text width=1.6cm, align=center, inner sep=2mm,
      font=\footnotesize\sffamily,
    },
    entradasalida/.style={
      draw=azulNoche, fill=verdeSignal!10!white,
      trapezium, trapezium left angle=75, trapezium right angle=105,
      text width=1.8cm, align=center, inner sep=2mm,
      font=\footnotesize\sffamily,
    },
    decision/.style={
      draw=azulNoche, fill=amarilloNota!20!white,
      diamond, aspect=1.6, text width=1.5cm, align=center, inner sep=1pt,
      font=\footnotesize\sffamily,
    },
    almacenamiento/.style={
      draw=azulNoche, fill=grisPapel,
      cylinder, shape border rotate=90, aspect=0.3,
      text width=1.7cm, align=center, inner sep=2mm,
      font=\footnotesize\sffamily,
    },
    documento/.style={
      draw=azulNoche, fill=azulMedio!8!white,
      text width=1.8cm, align=center, inner sep=2mm,
      font=\footnotesize\sffamily,
    },
    nota/.style={
      draw=grisLinea, dashed, fill=grisPapel,
      text width=1.8cm, align=center, inner sep=2mm,
      font=\scriptsize\sffamily,
    },
    etiqueta/.style={
      font=\scriptsize\sffamily, fill=white,
      fill opacity=0.75, text opacity=1, inner sep=1pt,
    },
  }


% ── Cabecera del documento ────────────────────────────────────────────────────
\fancyhead[L]{\small\color{azulNoche}\textbf{ELECTRONICA}}
\fancyhead[R]{\small\color{grisTexto}flujo\_motor\_fallback}
\renewcommand{\headrulewidth}{0pt}

% ─────────────────────────────────────────────────────────────────────────────
\begin{document}

\bandaTitulo{Failover del motor del orquestador}{Deteccion determinista del agotamiento de la cuota free y conmutacion al motor de respaldo con auto-restore}

\vspace{4pt}
\begin{center}
\begin{tcolorbox}[cajaContenido, title={\faIcon{map-signs}~Leyenda de símbolos---ISO 5807 / ANSI X3.5}]
\small
Óvalo = \textbf{terminador} (inicio/fin) $\cdot$
Rectángulo = \textbf{proceso} (acción determinista) $\cdot$
Rombo = \textbf{decisión} (ramas Sí/No) $\cdot$
Paralelogramo = \textbf{entrada/salida} (lectura/escritura de archivos) $\cdot$
Cilindro = \textbf{almacenamiento online} (estado, snapshots) $\cdot$
Rectángulo con base ondulada = \textbf{documento} (sesiones, logs) $\cdot$
Rectángulo punteado gris = \textbf{nota explicativa} (sin acción).
Flujo: arriba$\rightarrow$abajo, izquierda$\rightarrow$derecha.

\faIcon{tags}~Cada símbolo lleva una \textbf{etiqueta} (T/P/D/E/A/N + número, según su tipo);
su significado se expande en la tabla \emph{Leyenda de etiquetas} bajo cada diagrama.
\end{tcolorbox}
\end{center}

\vspace{6pt}

\section{\iconotexto{exchange-alt}{Deteccion de cuota agotada y conmutacion}}

\begin{center}
\begin{tikzpicture}[flujo, >=Stealth, x=1cm, y=1cm]
  % Fondo decorativo
  \fill[azulNoche!3!white, rounded corners=4pt]
    (-2.3,1.8) rectangle (9.3,-19.9);
  \node[terminador] (T1) at (0.0,0.0) {T1};
  \node[proceso] (P1) at (0.0,-2.2) {P1};
  \node[decision] (D1) at (0.0,-4.4) {D1};
  \node[terminador] (T3) at (0.0,-6.6) {T3};
  \node[decision] (D2) at (3.5,-4.4) {D2};
  \node[terminador] (T4) at (7.0,-6.6) {T4};
  \node[almacenamiento] (A1) at (3.5,-8.8) {A1};
  \node[proceso] (P2) at (3.5,-11.0) {P2};
  \node[proceso] (P3) at (3.5,-13.2) {P3};
  \node[proceso] (W1) at (3.5,-15.4) {W1};
  \node[decision] (D3) at (7.0,-15.4) {D3};
  \node[terminador] (T5) at (3.5,-17.6) {T5};
  \node[proceso] (P4) at (7.0,-17.6) {P4};

  \draw[->] (T1.south) -- (P1.north);
  \draw[->] (P1.south) -- (D1.north);
  \draw[->] (D1.south) -- (T3.north) node[etiqueta, left]{No};
  \draw[->] (D1.east) -- (D2.west) node[etiqueta, above]{Si};
  \draw[->] (D2.east) -- (T4.west) node[etiqueta, above]{No};
  \draw[->] (D2.south) -- (A1.north) node[etiqueta, left]{Si};
  \draw[->] (A1.south) -- (P2.north);
  \draw[->] (P2.south) -- (P3.north);
  \draw[->] (P3.south) -- (W1.north);
  \draw[->] (W1.east) -- (D3.west);
  \draw[->] (D3.west) -- (T5.east) node[etiqueta, above]{No};
  \draw[->] (D3.south) -- (P4.north) node[etiqueta, left]{Si};
\end{tikzpicture}
\end{center}

\vspace{6pt}
\begin{tcolorbox}[cajaContenido, title={\faIcon{table}~Leyenda de etiquetas}]
\small
\begin{tabularx}{\linewidth}{@{}>{\centering\arraybackslash}p{1.4cm}X@{}}
\toprule
\textbf{Etiqueta} & \textbf{Significado} \\ \midrule
A1 & Backup config previa + marker .tmp/motor\_fallback.json (timestamp, motivo, previo). \\
D1 & Firma de cuota agotada (429|quota|rate|limit|usage|reset|retry|retry\_after)? \\
D2 & Confirmaciones consecutivas >= N (retry budget max 3)? \\
D3 & Cuota free restablecida (2 sondeos ok consecutivos)? \\
P1 & Sonda: execution/verificar\_cuota\_motor.py check --modelo vigente. Clasifica: ok|agotada|error\_infra|indeterminado. \\
P2 & Switch: execution/aplicar\_switch\_modelo.py switch --modelo openrouter/qwen/qwen3.8-max-0902. \\
P3 & Alertar + relanzar TUI en tmux con el motor de respaldo. \\
P4 & Auto-restore: execution/aplicar\_switch\_modelo.py restore + relaunch TUI. \\
T1 & Inicio: motor free vigente (opencode/big-pickle) en la TUI. \\
T3 & No: cuota free OK, seguir con el motor vigente. \\
T4 & No: confirmacion insuficiente, re-sondear. \\
T5 & Seguir con el motor de respaldo; auto-restore pendiente. \\
W1 & Watcher daemon (--watch): re-sondear cuota free cada 900s. \\
\bottomrule
\end{tabularx}
\end{tcolorbox}


\section{\iconotexto{sticky-note}{Notas}}
\begin{itemize}
\item La conmutacion es funcion pura de la sonda (nunca razonada en el chat)
\item El motor de respaldo (openrouter/qwen/qwen3.8-max-0902) consume saldo OpenRouter; auto-restore la minimiza
\item Sonda activa: \textasciitilde{}8-9k tokens free; modo pasivo escanea opencode.log sin gastar tokens
\item Auto-restore: 2 sondeos ok consecutivos => restore + relaunch TUI con motor free
\end{itemize}

% ─────────────────────────────────────────────────────────────────────────────
\end{document}
